\documentclass[10pt,a4paper]{amsart}
\usepackage[margin=19mm]{geometry}
\usepackage{amsmath,amssymb,mathtools}
\usepackage[scr=rsfs,cal=euler]{mathalfa}
\usepackage{xfrac}
\usepackage[unicode,hidelinks,bookmarks=false]{hyperref}
\newcommand{\ie}{i.e.\ }
\newcommand{\cf}{cf.\ }
\newcommand{\resp}{respectively\ }
\newcommand{\Subsection}{\subsection}
\providecommand{\nobreakdash}{\nobreak\hbox{-}}
\setlength{\emergencystretch}{2em}
\multlinegap0pt
\usepackage{mathrsfs}
\usepackage{enumitem}
\usepackage{bbm}
\usepackage{lineno}
\newtheorem{Theorem}{Theorem}[section]
\newtheorem{Cor}[Theorem]{Corollary}
\newtheorem{Lemma}[Theorem]{Lemma}
\newtheorem{Proposition}[Theorem]{Proposition}
\newtheorem{Definition}[Theorem]{Definition}
\newtheorem{rem}[Theorem]{Remark}
\newtheorem{example}[Theorem]{Example}
\newtheorem{notation}[Theorem]{Notation}
\begin{document}
\title[Entropy Geometry and Normalized Means on Infinite-Dimensiona]{Entropy Geometry and Normalized Means on Infinite-Dimensional Hamiltonian Manifolds}
\author{Jean-Pierre Magnot}
\date{}
\maketitle
\noindent\emph{Source and licence.} Entropy Geometry and Normalized Means on Infinite-Dimensional Hamiltonian Manifolds. Version arXiv:2607.28660v1 (CC BY 4.0). This material is distributed under the Creative Commons Attribution 4.0 International licence, http://creativecommons.org/licenses/by/4.0/, as recorded in the arXiv OAI licence metadata for this exact version and shown on the abstract page https://arxiv.org/abs/2607.28660v1. Full rights notice: licenses/arxiv-2607.28660-CC-BY-4.0.txt.\par\smallskip
\noindent\textit{Editorial scope.} Counted unit: the original \texttt{Theorem} environment \texttt{thm:geom-eq} at source lines 961--1036 together with its complete proof at lines 1038--1084; the delivered document keeps source lines 292--1954 so that every referenced label and every citation resolves inside it. Native editable source is the paper's own concatenated TeX; the AMS wrapper is editorial formatting only. No publisher class, upstream script or unrelated artwork is redistributed.
\section{Infinite-Dimensional Hamiltonian Geometry and Normalized Means}

This section introduces the functional and geometric framework used
throughout the paper.

%------------------------------------------------------------
\subsection{Hamiltonian geometry on Fréchet manifolds}
%------------------------------------------------------------

Let $\mathcal M$ be a Fréchet manifold modelled on a locally convex
topological vector space $E$. We assume that $\mathcal M$ is equipped
with a smooth closed $2$-form $\omega$ such that, for every
$m\in\mathcal M$, the induced map
\[
\omega_m^\flat\colon T_m\mathcal M\longrightarrow T_m^*\mathcal M,
\qquad
v\longmapsto\omega_m(v,\cdot),
\]
is continuous and injective. Thus $\omega$ is a \emph{weak symplectic
form}. Writing
\[
J_m=\omega_m^\flat,
\]
we have
\[
\omega_m(v,w)=\langle J_m v,w\rangle.
\]
The pair $(\mathcal M,\omega)$ is called a \emph{weak symplectic
Fréchet manifold} (see \cite{Hamilton82,Neeb06,KM97}).

A \emph{Hamiltonian function} is a smooth map
\[
H\colon\mathcal M\longrightarrow\mathbb R
\]
such that
\begin{equation}
\omega\bigl(X_H,\cdot\bigr)=\mathrm dH
\label{eq:Hamiltonian}
\end{equation}
admits a smooth solution $X_H$. The associated vector field generates a
possibly local Hamiltonian flow $\Phi_t^H$.

Let $G$ be a Fréchet--Lie group acting smoothly on $\mathcal M$ by
symplectomorphisms, with Lie algebra $\mathfrak g$. A \emph{moment map}
is a smooth map
\[
\mathbf J\colon\mathcal M\longrightarrow\mathfrak g^*
\]
such that, for every $\xi\in\mathfrak g$, the function
\[
m\longmapsto\langle\mathbf J(m),\xi\rangle
\]
is a Hamiltonian generating the infinitesimal action of $\xi$. This
extends the classical finite-dimensional theory
\cite{AbrahamMarsden78,MarsdenRatiu99,Omori97}.

\begin{Definition}[Cylindrical functional]
A function $f\colon\mathcal M\to\mathbb R$ is called cylindrical if it
factors through a finite-dimensional smooth map, i.e.
\[
f=\widetilde f\circ\pi_N
\]
for some smooth map
\[
\pi_N\colon\mathcal M\longrightarrow\mathbb R^N
\]
and some function $\widetilde f$ on $\mathbb R^N$.
\end{Definition}

%------------------------------------------------------------
\subsection{Normalized means}
%------------------------------------------------------------

We now introduce the notion replacing probability measures in infinite
dimension.

Let
\[
\mathcal A_b\subset C_b(\mathcal M)
\]
be a unital algebra of bounded continuous observables, stable under the
exponential operations considered below. Let $\mathcal L$ be a unital
vector lattice algebra of real-valued functions on $\mathcal M$ such
that
\[
\mathcal A_b\subset\mathcal L.
\]
The space $\mathcal L$ may contain unbounded functions, including the
Hamiltonians and constraint observables used in the applications.

\begin{Definition}[Normalized mean]
\label{def:normalized-mean}
A \emph{normalized mean} on $\mathcal L$ is a linear functional
\[
\mathsf m\colon\mathcal L\longrightarrow\mathbb R
\]
such that:
\begin{enumerate}
\item[\textup{(i)}] $\mathsf m(1)=1$;
\item[\textup{(ii)}] $\mathsf m(f)\geq 0$ for every $f\in\mathcal L$
such that $f\geq 0$.
\end{enumerate}
\end{Definition}

We denote by $\mathsf S(\mathcal L)$ the set of all normalized means on
$\mathcal L$.

Normalized means encode finite additivity on every family of sets whose
indicator functions belong to $\mathcal L$, but they do not rely on
$\sigma$-additivity.

\begin{Definition}[Invariance]
A normalized mean $\mathsf m_0$ is said to be invariant under a flow
$\Phi_t$ if, for all $f\in\mathcal L$,
\[
\mathsf m_0(f\circ\Phi_t)=\mathsf m_0(f)
\]
whenever $f\circ\Phi_t\in\mathcal L$.
\end{Definition}

They generalize invariant means from amenable group theory.

\begin{rem}[Relation with amenability]
Invariant means on groups were introduced by von Neumann and further
developed by Day and Paterson
\cite{vonNeumann29,Day57,Paterson88}. The present framework extends this
concept to general configuration spaces.
\end{rem}

\begin{example}[Cut-off construction]
\label{ex:cutoff-mean}
Let $(B_N,\mu_N)$ be a sequence of finite-dimensional approximations of
$\mathcal M$, with
\[
0<\int_{B_N}\mathrm d\mu_N<\infty.
\]
If, for every $f\in\mathcal L$ under consideration, the limit
\[
\mathsf m_0(f)
=
\lim_{N\to\infty}
\frac{\displaystyle\int_{B_N}f\,\mathrm d\mu_N}
{\displaystyle\int_{B_N}\mathrm d\mu_N}
\]
exists, is finite, and is independent of the admissible truncation, then
$\mathsf m_0$ defines a normalized mean on its domain
\cite{Magnot17}.
\end{example}

%------------------------------------------------------------
\subsection{Topology on the space of normalized means}
\label{subsec:topology-means}
%------------------------------------------------------------

The definition of a normalized mean is algebraic: positivity and
normalization do not require continuity with respect to the compact-open
topology. This distinction is important for means obtained from
finite-dimensional cut-offs, which need not be continuous for that
topology.

We equip $\mathsf S(\mathcal L)$ with the topology of pointwise
convergence on $\mathcal L$. Thus, a net
$(\mathsf m_i)_{i\in I}$ converges to $\mathsf m$ if and only if
\[
\mathsf m_i(f)\longrightarrow\mathsf m(f)
\qquad
\text{for every }f\in\mathcal L.
\]
Equivalently, this is the weakest topology for which all evaluation maps
\[
\operatorname{ev}_f\colon
\mathsf S(\mathcal L)\longrightarrow\mathbb R,
\qquad
\operatorname{ev}_f(\mathsf m)=\mathsf m(f),
\]
are continuous.

When $\mathcal A_b\subset C_b(\mathcal M)$ is the algebra of bounded
test observables, positivity and normalization imply
\[
\left|\mathsf m(f)\right|
\leq
\|f\|_\infty,
\qquad
f\in\mathcal A_b,
\quad
\mathsf m\in\mathsf S(\mathcal A_b).
\]
Indeed,
\[
-\|f\|_\infty 1
\leq
f
\leq
\|f\|_\infty 1,
\]
and positivity gives the required estimate.

Consequently, $\mathsf S(\mathcal A_b)$ may be identified with a closed
subset of
\[
\prod_{f\in\mathcal A_b}
[-\|f\|_\infty,\|f\|_\infty].
\]
By Tychonoff's theorem, it is compact for the topology of pointwise
convergence on $\mathcal A_b$.

\begin{Proposition}
\label{prop:compact-bounded-states}
The space $\mathsf S(\mathcal A_b)$ of normalized means on the bounded
observable algebra $\mathcal A_b$ is compact for the topology of
pointwise convergence on $\mathcal A_b$.
\end{Proposition}

\begin{proof}
For every $f\in\mathcal A_b$, positivity and normalization give
\[
|\mathsf m(f)|\leq\|f\|_\infty.
\]
Hence $\mathsf S(\mathcal A_b)$ is contained in the compact product
\[
K
=
\prod_{f\in\mathcal A_b}
[-\|f\|_\infty,\|f\|_\infty].
\]
Linearity, positivity, and normalization are preserved under pointwise
limits. Therefore $\mathsf S(\mathcal A_b)$ is closed in $K$, and is
thus compact.
\end{proof}

For the larger algebra $\mathcal L$, which may contain unbounded
observables, no analogous compactness statement holds in general.
Accordingly, compactness of the entire space $\mathsf S(\mathcal L)$
will not be assumed. Variational arguments based on compactness require
instead compactness of the relevant free-energy sublevel sets,
\[
\left\{
\mathsf n\in\mathsf S(\mathcal L)
\;\middle|\;
\mathcal G_{\beta,\lambda}(\mathsf n)\leq r
\right\},
\]
for the topology of pointwise convergence. Whenever such an argument is
used, this property will be included explicitly among the coercivity
assumptions.

%------------------------------------------------------------
\subsection{Pointwise limits and cut-off normalized means}
\label{subsec:cutoff-limits}
%------------------------------------------------------------

The topology of pointwise convergence is naturally adapted to
normalized means constructed from finite-dimensional approximations.

Let $(B_N,\mu_N)$ be a sequence of finite-dimensional approximations
such that
\[
0<\int_{B_N}\mathrm d\mu_N<\infty,
\]
and define
\[
\mathsf m_N(f)
=
\frac{
\displaystyle\int_{B_N}f\,\mathrm d\mu_N
}{
\displaystyle\int_{B_N}\mathrm d\mu_N
}
\]
whenever the quotient is well defined.

\begin{Proposition}
\label{prop:cutoff-pointwise-limit}
Assume that, for every $f\in\mathcal A_b$, the limit
\[
\mathsf m_0(f)
=
\lim_{N\to\infty}\mathsf m_N(f)
\]
exists. Then $\mathsf m_0$ is a normalized mean on $\mathcal A_b$.

More generally, if the same pointwise limit exists and is finite for
every $f$ in a vector lattice algebra $\mathcal L$, then $\mathsf m_0$
is a normalized mean on $\mathcal L$.
\end{Proposition}

\begin{proof}
For every $N$, the functional $\mathsf m_N$ is linear, positive, and
normalized. These three properties are preserved under pointwise
limits. In particular,
\[
\mathsf m_0(1)
=
\lim_{N\to\infty}\mathsf m_N(1)
=
1,
\]
and, whenever $f\geq 0$,
\[
\mathsf m_0(f)
=
\lim_{N\to\infty}\mathsf m_N(f)
\geq 0.
\]
Therefore $\mathsf m_0$ is a normalized mean.
\end{proof}

\begin{rem}
No continuity for the compact-open topology is required in
Proposition~\ref{prop:cutoff-pointwise-limit}. In fact, means describing
asymptotic behavior at infinity are generally not continuous for that
topology.
\end{rem}

\begin{rem}
The preceding proposition proves the existence of a normalized mean
once pointwise convergence has been established. It does not by itself
prove that the resulting mean is independent of the chosen
finite-dimensional exhaustion. Such independence must either be
verified in each application or included in the definition of an
admissible family of cut-offs.
\end{rem}

%------------------------------------------------------------
\subsection{Exponentially admissible potentials}
%------------------------------------------------------------

The Hamiltonians arising in the applications are generally unbounded.
They are therefore treated as thermodynamic potentials rather than as
elements of the algebra $\mathcal A_b$ of bounded test observables.

\begin{Definition}[Exponentially admissible potential]
\label{def:admissible-potential}
A real-valued function $F\in\mathcal L$ is called
\emph{$\mathsf m_0$-admissible} if
\[
e^F\in\mathcal L,
\qquad
0<\mathsf m_0(e^F)<\infty,
\]
and
\[
fe^F\in\mathcal L
\qquad
\text{for every }f\in\mathcal L.
\]
The set of $\mathsf m_0$-admissible potentials is denoted by
$\mathcal E(\mathsf m_0)$.
\end{Definition}

We assume that $\mathcal E(\mathsf m_0)$ contains $0$ and is stable
under the bounded perturbations used below. For
$F\in\mathcal E(\mathsf m_0)$, define
\[
\Lambda_{\mathsf m_0}(F)
=
\log\mathsf m_0(e^F).
\]

\begin{Definition}[Exponential tilt]
\label{def:exponential-tilt}
Let $F\in\mathcal E(\mathsf m_0)$. The exponential tilt of
$\mathsf m_0$ by $F$ is the normalized mean
\[
\mathsf m_F(f)
=
\frac{\mathsf m_0(fe^F)}{\mathsf m_0(e^F)},
\qquad
f\in\mathcal L.
\]
\end{Definition}

The positivity and normalization of $\mathsf m_F$ follow immediately:
\[
\mathsf m_F(f)\geq 0
\quad\text{whenever }f\geq 0,
\qquad
\mathsf m_F(1)=1.
\]

%------------------------------------------------------------
\subsection{Relative entropy}
%------------------------------------------------------------

Given $\mathsf m_0\in\mathsf S(\mathcal L)$, we define entropy through a
variational principle.

\begin{Definition}[Relative entropy]
\label{def:entropy}
For $\mathsf n,\mathsf m_0\in\mathsf S(\mathcal L)$, define
\[
\mathcal H(\mathsf n\,\|\,\mathsf m_0)
=
\sup_{f\in\mathcal E(\mathsf m_0)}
\left\{
\mathsf n(f)-\log\mathsf m_0(e^f)
\right\},
\]
with values in $[0,+\infty]$.
\end{Definition}

This definition is inspired by the Donsker--Varadhan variational
formula.

\begin{Proposition}[Basic properties of the relative entropy]
\label{prop:entropy-properties}
Let $\mathsf m_0$ be a normalized mean on $\mathcal L$. Assume that
there exists a vector subspace $\mathcal T\subset\mathcal L$ separating
normalized means such that, for every $h\in\mathcal T$, one has
$th\in\mathcal E(\mathsf m_0)$ for all sufficiently small
$t\in\mathbb R$, and
\[
\left.\frac{\mathrm d}{\mathrm dt}\right|_{t=0}
\log\mathsf m_0(e^{th})
=
\mathsf m_0(h).
\]
Then:
\begin{enumerate}
\item
\[
\mathcal H(\mathsf n\,\|\,\mathsf m_0)\geq 0,
\]
with equality if and only if $\mathsf n=\mathsf m_0$;
\item the map
\[
\mathsf n\longmapsto
\mathcal H(\mathsf n\,\|\,\mathsf m_0)
\]
is convex and lower semicontinuous for the topology of pointwise
convergence on $\mathsf S(\mathcal L)$.
\end{enumerate}
\end{Proposition}

\begin{proof}
The non-negativity follows by taking $f=0$. Convexity and lower
semicontinuity follow because
\[
\mathsf n\longmapsto
\mathsf n(f)-\log\mathsf m_0(e^f)
\]
is affine and continuous for the topology of pointwise convergence, for
every admissible $f$.

Suppose now that
\[
\mathcal H(\mathsf n\,\|\,\mathsf m_0)=0.
\]
For every $h\in\mathcal T$ and every sufficiently small
$t\in\mathbb R$, one has
\[
t\mathsf n(h)
\leq
\log\mathsf m_0(e^{th}).
\]
Dividing by $t$ and letting $t\to0$ from the positive and negative sides
gives
\[
\mathsf n(h)=\mathsf m_0(h).
\]
Since $\mathcal T$ separates normalized means, it follows that
$\mathsf n=\mathsf m_0$.
\end{proof}

%------------------------------------------------------------
\subsection{Change-of-reference compatibility}
%------------------------------------------------------------

The variational entropy is defined through a reference-dependent class
of admissible potentials. Consequently, exponential tilting requires a
compatibility condition between the admissible classes associated with
the original and tilted means.

\begin{Definition}[Change-of-reference compatibility]
\label{def:change-reference-compatible}
Let $F\in\mathcal E(\mathsf m_0)$. We say that the exponential tilt by
$F$ is \emph{change-of-reference compatible} if
\[
\mathcal H(\mathsf n\,\|\,\mathsf m_F)
=
\mathcal H(\mathsf n\,\|\,\mathsf m_0)
-\mathsf n(F)
+\log\mathsf m_0(e^F)
\]
for every normalized mean $\mathsf n$ for which the terms are defined.
\end{Definition}

A sufficient translation condition on the admissible classes, together
with a proof of the resulting change-of-reference identity, is given in
section~\ref{subsec:change-reference}.

%------------------------------------------------------------
\subsection{Analytic assumptions}
%------------------------------------------------------------

Let $H,\Phi_1,\ldots,\Phi_N\in\mathcal L$ be possibly unbounded
observables, and set
\[
F_{\beta,\lambda}
=
-\beta H-\sum_{a=1}^{N}\lambda^a\Phi_a.
\]
The admissible thermodynamic domain is
\[
\mathcal D
=
\left\{
(\beta,\lambda)\in\mathbb R^{N+1}
\;\middle|\;
F_{\beta,\lambda}\in\mathcal E(\mathsf m_0)
\right\}.
\]

Depending on the result under consideration, we shall use the following
assumptions:

\begin{itemize}
\item[(A1)] Exponential admissibility: the relevant parameters belong
to $\mathcal D$, so that
\[
0<
\mathsf m_0\left(
e^{-\beta H-\sum_{a=1}^{N}\lambda^a\Phi_a}
\right)
<\infty;
\]

\item[(A2)] Coercivity: the sublevel sets of
$\mathcal G_{\beta,\lambda}$ are compact in
$\mathsf S(\mathcal L)$ for the topology of pointwise convergence;

\item[(A3)] Non-degeneracy: the covariance form has no nonzero null
direction on the reduced parameter domain under consideration.
\end{itemize}

%------------------------------------------------------------
\subsection{Variational principle}
%------------------------------------------------------------

For $(\beta,\lambda)\in\mathcal D$, define
\[
Z(\beta,\lambda)
=
\mathsf m_0\left(e^{F_{\beta,\lambda}}\right).
\]

The free-energy functional associated with $(\beta,\lambda)$ is
\[
\mathcal G_{\beta,\lambda}(\mathsf n)
=
\mathcal H(\mathsf n\,\|\,\mathsf m_0)
+
\beta\mathsf n(H)
+
\sum_{a=1}^{N}\lambda^a\mathsf n(\Phi_a),
\qquad
\mathsf n\in\mathsf S(\mathcal L).
\]
Equivalently,
\[
\mathcal G_{\beta,\lambda}(\mathsf n)
=
\mathcal H(\mathsf n\,\|\,\mathsf m_0)
-
\mathsf n(F_{\beta,\lambda}).
\]

\begin{Proposition}[Gibbs variational principle]
\label{prop:gibbs-variational-principle}
Let $(\beta,\lambda)\in\mathcal D$ and assume that the exponential tilt
by $F_{\beta,\lambda}$ is change-of-reference compatible. Assume also
that the relative entropy with reference
$\mathsf n_{\beta,\lambda}^*$ has the separation property of
Proposition~\ref{prop:entropy-properties}. Then
\[
-\log Z(\beta,\lambda)
=
\inf_{\mathsf n\in\mathsf S(\mathcal L)}
\mathcal G_{\beta,\lambda}(\mathsf n).
\]
Equivalently,
\[
\log Z(\beta,\lambda)
=
\sup_{\mathsf n\in\mathsf S(\mathcal L)}
\left\{
\mathsf n(F_{\beta,\lambda})
-
\mathcal H(\mathsf n\,\|\,\mathsf m_0)
\right\}.
\]
The extremum is attained at the unique normalized mean
\[
\mathsf n_{\beta,\lambda}^{*}(f)
=
\frac{
\mathsf m_0\left(fe^{F_{\beta,\lambda}}\right)
}{
\mathsf m_0\left(e^{F_{\beta,\lambda}}\right)
},
\qquad
f\in\mathcal L.
\]
\end{Proposition}

\begin{proof}
By the change-of-reference identity,
\[
\mathcal H\left(
\mathsf n\,\|\,\mathsf n_{\beta,\lambda}^{*}
\right)
=
\mathcal H(\mathsf n\,\|\,\mathsf m_0)
-\mathsf n(F_{\beta,\lambda})
+\log Z(\beta,\lambda).
\]
Therefore
\[
\mathcal H\left(
\mathsf n\,\|\,\mathsf n_{\beta,\lambda}^{*}
\right)
=
\mathcal G_{\beta,\lambda}(\mathsf n)
+\log Z(\beta,\lambda).
\]
Since relative entropy is nonnegative,
\[
\mathcal G_{\beta,\lambda}(\mathsf n)
\geq
-\log Z(\beta,\lambda).
\]
For
\[
\mathsf n=\mathsf n_{\beta,\lambda}^{*},
\]
the relative entropy on the left-hand side vanishes, and hence equality
holds. If equality holds for another normalized mean $\mathsf n$, then
\[
\mathcal H\left(
\mathsf n\,\|\,\mathsf n_{\beta,\lambda}^{*}
\right)=0.
\]
The separation property of the relative entropy with reference
$\mathsf n_{\beta,\lambda}^{*}$ therefore implies
\[
\mathsf n=\mathsf n_{\beta,\lambda}^{*}.
\]
\end{proof}

\begin{rem}
The compactness assumption \textup{(A2)} is not needed in the preceding
proposition because the minimizer is explicitly constructed by
exponential tilting. It becomes relevant for existence arguments in
which no compatible exponential tilt is assumed in advance.
\end{rem}

%------------------------------------------------------------
\subsection{Entropy geometry}
%------------------------------------------------------------

We now relate the previous construction to symplectic geometry.

For a left action of $G$ on $\mathcal M$, we use the induced action on
observables
\[
(g\cdot f)(m)=f(g^{-1}\cdot m).
\]


\begin{Theorem}[Geometric equilibrium for normalized means]
\label{thm:geom-eq}
Let $(\mathcal M,\omega,G,\mathbf J)$ be a weak Hamiltonian Fréchet
$G$-space, and let $\mathsf m_0$ be a $G$-invariant normalized mean on
$\mathcal L$. Assume that the action of $G$ preserves $\mathcal L$.

Assume that $H\in\mathcal L$ is $G$-invariant and that the moment map
\[
\mathbf J\colon\mathcal M\longrightarrow\mathfrak g^*
\]
is equivariant. For $\lambda\in\mathfrak g$ and
$(\beta,\lambda)$ in the admissible domain, define
\[
F_{\beta,\lambda}
=
-\beta H-\langle\mathbf J,\lambda\rangle
\]
and
\[
Z(\beta,\lambda)
=
\mathsf m_0\left(e^{F_{\beta,\lambda}}\right).
\]
The corresponding equilibrium mean is
\[
\mathsf n_{\beta,\lambda}^{*}(f)
=
\frac{
\mathsf m_0\left(fe^{F_{\beta,\lambda}}\right)
}{
\mathsf m_0\left(e^{F_{\beta,\lambda}}\right)
},
\qquad
f\in\mathcal L.
\]

Then the family of equilibrium means is equivariant with respect to the
adjoint action on the parameter:
\[
\mathsf n_{\beta,\lambda}^{*}(g\cdot f)
=
\mathsf n_{\beta,\operatorname{Ad}_{g^{-1}}\lambda}^{*}(f),
\qquad
g\in G.
\]
In particular, $\mathsf n_{\beta,\lambda}^{*}$ is invariant under the
stabilizer
\[
G_\lambda
=
\left\{
g\in G
\;\middle|\;
\operatorname{Ad}_{g^{-1}}\lambda=\lambda
\right\}.
\]
If $\lambda$ is fixed by the adjoint action, then
$\mathsf n_{\beta,\lambda}^{*}$ is $G$-invariant.

Moreover, the thermodynamic potential
\[
\psi(\beta,\lambda)=\log Z(\beta,\lambda)
\]
generates the correspondence between the intensive variables
$(\beta,\lambda)$ and the extensive variables
\[
E(\beta,\lambda)
=
\mathsf n_{\beta,\lambda}^{*}(H),
\qquad
c(\beta,\lambda)
=
\mathsf n_{\beta,\lambda}^{*}(\mathbf J),
\]
on every domain on which $\psi$ is differentiable.
\end{Theorem}

\begin{proof}
Equivariance of the moment map means
\[
\mathbf J(g\cdot m)=\operatorname{Ad}_g^*\mathbf J(m).
\]
Consequently,
\[
F_{\beta,\lambda}(g\cdot m)
=
-\beta H(m)
-
\left\langle
\operatorname{Ad}_g^*\mathbf J(m),\lambda
\right\rangle
=
F_{\beta,\operatorname{Ad}_{g^{-1}}\lambda}(m).
\]
Using the $G$-invariance of $\mathsf m_0$, we obtain
\[
\begin{split}
\mathsf n_{\beta,\lambda}^{*}(g\cdot f)
&=
\frac{
\mathsf m_0\left((g\cdot f)e^{F_{\beta,\lambda}}\right)
}{
\mathsf m_0\left(e^{F_{\beta,\lambda}}\right)
}
\\
&=
\frac{
\mathsf m_0\left(
f e^{F_{\beta,\operatorname{Ad}_{g^{-1}}\lambda}}
\right)
}{
\mathsf m_0\left(
e^{F_{\beta,\operatorname{Ad}_{g^{-1}}\lambda}}
\right)
}
\\
&=
\mathsf n_{\beta,\operatorname{Ad}_{g^{-1}}\lambda}^{*}(f).
\end{split}
\]
The assertions concerning the stabilizer follow immediately. The
correspondence between intensive and extensive variables follows from
the first-derivative formulas established below.
\end{proof}


\begin{rem}
The preceding theorem establishes equivariance, and in particular
invariance under the stabilizer of the thermodynamic parameter. A
KMS-type property requires additional Poisson-algebraic assumptions and
does not follow from invariance alone.
\end{rem}

%--------------------------------------------------------------
\subsection{Local exponential regularity}
\label{subsec:local-exponential-regularity}
%--------------------------------------------------------------

Pointwise finiteness of the partition functional does not by itself
imply differentiability with respect to the thermodynamic parameters.
We therefore introduce a local exponential moment condition.

Let
\[
\theta=(\beta,\lambda^1,\ldots,\lambda^N)\in\mathcal D
\]
and set
\[
\mathcal O_0=H,
\qquad
\mathcal O_a=\Phi_a,
\quad
a=1,\ldots,N.
\]

\begin{Definition}[Local exponential regularity]
\label{def:local-exponential-regularity}
The partition functional is said to be
\emph{locally exponentially regular at $\theta\in\mathcal D$} if
$\theta$ is an interior point of $\mathcal D$ and there exists
$\varepsilon>0$ such that
\[
\exp\left(
F_\theta+\varepsilon\sum_{j=0}^{N}|\mathcal O_j|
\right)
\in\mathcal L
\]
and
\[
\mathsf m_0\left(
\exp\left(
F_\theta+\varepsilon\sum_{j=0}^{N}|\mathcal O_j|
\right)
\right)
<\infty.
\]
We say that the partition functional is locally exponentially regular
on $\mathcal D$ if this condition holds at every
$\theta\in\mathcal D$.
\end{Definition}

The local exponential regularity condition provides a common
exponential bound for sufficiently small parameter increments. Indeed,
if
\[
\delta\theta
=
(\delta\beta,\delta\lambda^1,\ldots,\delta\lambda^N)
\]
satisfies
\[
|\delta\beta|<\varepsilon,
\qquad
|\delta\lambda^a|<\varepsilon,
\quad
a=1,\ldots,N,
\]
then
\[
F_{\theta+\delta\theta}
\leq
F_\theta
+
\varepsilon\sum_{j=0}^{N}|\mathcal O_j|.
\]

\begin{Proposition}[Smoothness of the partition functional]
\label{prop:partition-smoothness}
Assume that the partition functional is locally exponentially regular
on $\mathcal D$. Then $\mathcal D$ is open and
\[
Z(\theta)=\mathsf m_0(e^{F_\theta})
\]
is smooth on $\mathcal D$. Since $Z(\theta)>0$, the logarithmic
partition functional
\[
\psi(\theta)=\log Z(\theta)
\]
is also smooth.

For every parameter direction $\dot\theta$, one has
\[
DZ(\theta)[\dot\theta]
=
-\mathsf m_0\left(
\mathcal O_{\dot\theta}e^{F_\theta}
\right),
\]
where
\[
\mathcal O_{\dot\theta}
=
\dot\beta H+\sum_{a=1}^{N}\dot\lambda^a\Phi_a.
\]
Consequently,
\[
D\psi(\theta)[\dot\theta]
=
-\mathsf n_\theta^*(\mathcal O_{\dot\theta}).
\]

More generally, for parameter directions
$\dot\theta_1,\ldots,\dot\theta_k$,
\[
D^kZ(\theta)
[\dot\theta_1,\ldots,\dot\theta_k]
=
(-1)^k
\mathsf m_0\left(
\mathcal O_{\dot\theta_1}\cdots
\mathcal O_{\dot\theta_k}
e^{F_\theta}
\right).
\]
The higher derivatives of $\psi$ are the corresponding joint cumulants
with respect to $\mathsf n_\theta^*$.
\end{Proposition}

\begin{proof}
The openness of $\mathcal D$ follows from
Definition~\ref{def:local-exponential-regularity}. Fix
$\theta\in\mathcal D$. The local exponential regularity condition
provides a common exponential bound for the parameter increments.

For a direction $\dot\theta$ and sufficiently small $t$, Taylor's
formula gives
\[
e^{F_{\theta+t\dot\theta}}
=
e^{F_\theta}
\left(
1-t\mathcal O_{\dot\theta}
+
R_2(t)
\right),
\]
where
\[
|R_2(t)|
\leq
\frac{t^2}{2}
\mathcal O_{\dot\theta}^{\,2}
e^{|t|\,|\mathcal O_{\dot\theta}|}.
\]
After decreasing the parameter neighborhood if necessary, the
right-hand side is bounded by a fixed element of $\mathcal L$ having a
finite $\mathsf m_0$-mean. Positivity of $\mathsf m_0$ and the explicit
factor $t^2$ therefore imply
\[
DZ(\theta)[\dot\theta]
=
-\mathsf m_0\left(
\mathcal O_{\dot\theta}e^{F_\theta}
\right).
\]

The same argument, applied to higher-order Taylor remainders, gives
\[
D^kZ(\theta)
[\dot\theta_1,\ldots,\dot\theta_k]
=
(-1)^k
\mathsf m_0\left(
\mathcal O_{\dot\theta_1}\cdots
\mathcal O_{\dot\theta_k}
e^{F_\theta}
\right).
\]
The common local exponential bound also gives continuity of these
derivatives. Hence $Z$ is smooth. Since $Z>0$ on $\mathcal D$, the same
holds for $\psi=\log Z$.
\end{proof}

\begin{rem}
No interchange of an uncontrolled infinite-dimensional limit with
differentiation is used here. The derivatives are justified directly
by finite-order Taylor estimates and the local exponential bound.
\end{rem}

%--------------------------------------------------------------
\subsection{Covariance Hessian and strict convexity}
\label{subsec:covariance-hessian}
%--------------------------------------------------------------

For
\[
\theta=(\beta,\lambda^1,\ldots,\lambda^N)\in\mathcal D,
\]
write
\[
F_\theta
=
-\beta H-\sum_{a=1}^{N}\lambda^a\Phi_a,
\qquad
\psi(\theta)
=
\log Z(\theta).
\]
For a parameter direction
\[
\dot\theta
=
(\dot\beta,\dot\lambda^1,\ldots,\dot\lambda^N),
\]
set
\[
\mathcal O_{\dot\theta}
=
\dot\beta H+\sum_{a=1}^{N}\dot\lambda^a\Phi_a.
\]
Then
\[
F_{\theta+t\dot\theta}
=
F_\theta-t\mathcal O_{\dot\theta}.
\]

Assume that the local exponential regularity hypothesis of
Definition~\ref{def:local-exponential-regularity} holds. Differentiation
under the normalized mean gives
\[
D\psi(\theta)[\dot\theta]
=
-\mathsf n_\theta^*(\mathcal O_{\dot\theta}).
\]
For two parameter directions $\dot\theta_1$ and $\dot\theta_2$, a second
differentiation gives
\[
D^2\psi(\theta)
[\dot\theta_1,\dot\theta_2]
=
\operatorname{Cov}_{\mathsf n_\theta^*}
\left(
\mathcal O_{\dot\theta_1},
\mathcal O_{\dot\theta_2}
\right),
\]
where
\[
\operatorname{Cov}_{\mathsf n}(f,g)
=
\mathsf n(fg)-\mathsf n(f)\mathsf n(g).
\]
In particular,
\[
D^2\psi(\theta)[\dot\theta,\dot\theta]
=
\operatorname{Var}_{\mathsf n_\theta^*}
(\mathcal O_{\dot\theta})
\geq 0,
\]
where
\[
\operatorname{Var}_{\mathsf n}(f)
=
\mathsf n(f^2)-\mathsf n(f)^2.
\]

\begin{Definition}[Thermodynamically null direction]
\label{def:null-direction}
A parameter direction $\dot\theta$ is called
\emph{thermodynamically null at $\theta$} if
\[
\operatorname{Var}_{\mathsf n_\theta^*}
(\mathcal O_{\dot\theta})=0.
\]
\end{Definition}

Every direction for which $\mathcal O_{\dot\theta}$ is constant is
thermodynamically null. If the equilibrium mean is not faithful, a
nonconstant observable may also have zero variance.

\begin{Proposition}[Strict-convexity criterion]
\label{prop:strict-convexity}
Let $U\subset\mathcal D$ be convex, and assume that $\psi$ is twice
differentiable on $U$. Then $\psi$ is convex on $U$.

If
\[
\operatorname{Var}_{\mathsf n_\theta^*}
(\mathcal O_{\dot\theta})>0
\]
for every $\theta\in U$ and every nonzero parameter direction
$\dot\theta$, then $\psi$ is strictly convex on $U$.

At every $\theta\in U$, the Hessian induces a positive-definite
quadratic form on
\[
T_\theta U\big/\ker D^2\psi(\theta).
\]
If these kernels have locally constant rank, the resulting quotient
spaces form a smooth vector bundle equipped with the metric induced by
the covariance form.
\end{Proposition}

\begin{proof}
The Hessian of $\psi$ is the covariance form. Hence
\[
D^2\psi(\theta)[\dot\theta,\dot\theta]
=
\operatorname{Var}_{\mathsf n_\theta^*}
(\mathcal O_{\dot\theta})
\geq 0,
\]
which proves convexity. Under the stated non-degeneracy assumption, the
restriction of $\psi$ to every nonconstant affine line contained in
$U$ has strictly positive second derivative and is therefore strictly
convex.

At a fixed $\theta$, the kernel of the positive semidefinite covariance
form is precisely $\ker D^2\psi(\theta)$. The induced quadratic form on
the quotient by this kernel is positive definite. If the kernels have
locally constant rank, they form a smooth subbundle, and the quotient
bundle inherits the corresponding positive-definite metric.
\end{proof}

\begin{rem}
The map
\[
f\longmapsto\log\mathsf m_0(e^f)
\]
cannot be strictly convex in the direction of constant functions,
because
\[
\log\mathsf m_0(e^{f+c})
=
c+\log\mathsf m_0(e^f).
\]
Strict convexity must therefore be understood modulo constants, or more
generally modulo covariance-null directions that persist along the
parameter domain.
\end{rem}

\begin{rem}
If $\mathsf n_\theta^*$ is faithful and
\[
\operatorname{Var}_{\mathsf n_\theta^*}(f)=0,
\]
then
\[
\mathsf n_\theta^*
\left(
\bigl(f-\mathsf n_\theta^*(f)\bigr)^2
\right)
=
0
\]
implies
\[
f=\mathsf n_\theta^*(f)
\]
in $\mathcal L$. In that case, the thermodynamically null directions
are exactly those for which
\[
\dot\beta H+\sum_{a=1}^{N}\dot\lambda^a\Phi_a
\]
is constant.
\end{rem}

%--------------------------------------------------------------
\subsection{Legendre--Fenchel duality and thermodynamic entropy}
\label{subsec:legendre-duality}
%--------------------------------------------------------------

Let
\[
\theta=(\beta,\lambda^1,\ldots,\lambda^N)
\]
and let
\[
x=(E,c_1,\ldots,c_N)
\]
denote the corresponding vector of extensive variables. We use the
pairing
\[
\langle\theta,x\rangle
=
\beta E+\sum_{a=1}^{N}\lambda^a c_a.
\]

The logarithmic partition functional is
\[
\psi(\theta)
=
\log\mathsf m_0
\left(
e^{-\beta H-\sum_{a=1}^{N}\lambda^a\Phi_a}
\right).
\]
Under the local exponential regularity hypothesis,
\[
\partial_\beta\psi(\theta)
=
-\mathsf n_\theta^*(H)
\]
and
\[
\partial_{\lambda^a}\psi(\theta)
=
-\mathsf n_\theta^*(\Phi_a).
\]
Therefore,
\[
\nabla\psi(\theta)
=
-x(\theta),
\]
where
\[
x(\theta)
=
\left(
\mathsf n_\theta^*(H),
\mathsf n_\theta^*(\Phi_1),
\ldots,
\mathsf n_\theta^*(\Phi_N)
\right).
\]

\begin{Definition}[Thermodynamic entropy]
\label{def:thermodynamic-entropy}
The thermodynamic entropy associated with $\psi$ is defined by
\[
\mathcal S(x)
=
\inf_{\theta\in\mathcal D}
\left\{
\psi(\theta)+\langle\theta,x\rangle
\right\}.
\]
Equivalently,
\[
\mathcal S(E,c)
=
\inf_{(\beta,\lambda)\in\mathcal D}
\left\{
\log Z(\beta,\lambda)
+
\beta E
+
\sum_{a=1}^{N}\lambda^a c_a
\right\}.
\]
\end{Definition}

\begin{Proposition}[Concavity and Legendre correspondence]
\label{prop:legendre}
Assume that $\mathcal D$ is convex and that $\psi$ is differentiable
and convex on $\mathcal D$. Then $\mathcal S$ is concave on its
effective domain.

For $\theta\in\mathcal D$, set
\[
x(\theta)=-\nabla\psi(\theta).
\]
Then $\theta$ realizes the infimum in
Definition~\ref{def:thermodynamic-entropy}, and
\[
\mathcal S(x(\theta))
=
\psi(\theta)+\langle\theta,x(\theta)\rangle.
\]

If $\psi$ descends to a strictly convex differentiable function on a
quotient of the convex parameter domain by a fixed space of
thermodynamically null directions, then the map
\[
\theta\longmapsto x(\theta)
\]
is injective on that quotient.
\end{Proposition}

\begin{proof}
For every fixed $\theta\in\mathcal D$, the function
\[
x\longmapsto
\psi(\theta)+\langle\theta,x\rangle
\]
is affine. Since an infimum of affine functions is concave,
$\mathcal S$ is concave.

Let
\[
x(\theta)=-\nabla\psi(\theta).
\]
For every $\theta'\in\mathcal D$, convexity of $\psi$ gives
\[
\psi(\theta')
\geq
\psi(\theta)
+
\left\langle
\nabla\psi(\theta),\theta'-\theta
\right\rangle.
\]
Using
\[
\nabla\psi(\theta)=-x(\theta),
\]
we obtain
\[
\psi(\theta')
+
\langle\theta',x(\theta)\rangle
\geq
\psi(\theta)
+
\langle\theta,x(\theta)\rangle.
\]
Thus $\theta$ realizes the infimum defining
$\mathcal S(x(\theta))$.

Finally, the gradient of a differentiable strictly convex function on a
convex domain is injective. Applying this statement to the induced
function on the quotient proves the last assertion.
\end{proof}

\begin{Proposition}[Constrained variational entropy]
\label{prop:constrained-entropy}
For an extensive vector
\[
x=(E,c_1,\ldots,c_N),
\]
define
\[
\mathcal S_{\mathrm{var}}(x)
=
\sup
\left\{
-\mathcal H(\mathsf n\,\|\,\mathsf m_0)
\;\middle|\;
\mathsf n\in\mathsf S(\mathcal L),
\quad
\mathsf n(H)=E,
\quad
\mathsf n(\Phi_a)=c_a
\right\}.
\]
Then
\[
\mathcal S_{\mathrm{var}}(x)
\leq
\mathcal S(x).
\]

Suppose that the hypotheses of
Proposition~\ref{prop:gibbs-variational-principle} hold. If
$x=x(\theta)$ for some $\theta\in\mathcal D$, then equality holds and
the supremum is attained at $\mathsf n_\theta^*$:
\[
\mathcal S_{\mathrm{var}}(x(\theta))
=
\mathcal S(x(\theta))
=
-\mathcal H(\mathsf n_\theta^*\,\|\,\mathsf m_0).
\]
\end{Proposition}

\begin{proof}
Let $\mathsf n$ satisfy the constraints associated with $x$. The
definition of relative entropy, applied to $F_\theta$, gives, for every
$\theta\in\mathcal D$,
\[
\mathcal H(\mathsf n\,\|\,\mathsf m_0)
\geq
\mathsf n(F_\theta)-\psi(\theta).
\]
Since
\[
\mathsf n(F_\theta)
=
-\langle\theta,x\rangle,
\]
we obtain
\[
-\mathcal H(\mathsf n\,\|\,\mathsf m_0)
\leq
\psi(\theta)+\langle\theta,x\rangle.
\]
Taking first the supremum over all states satisfying the constraints
and then the infimum over $\theta$ gives
\[
\mathcal S_{\mathrm{var}}(x)
\leq
\mathcal S(x).
\]

If $x=x(\theta)$, the equilibrium mean $\mathsf n_\theta^*$ satisfies
the required constraints. By the change-of-reference identity,
evaluated at $\mathsf n_\theta^*$,
\[
\mathcal H(\mathsf n_\theta^*\,\|\,\mathsf m_0)
=
\mathsf n_\theta^*(F_\theta)-\psi(\theta).
\]
Since
\[
\mathsf n_\theta^*(F_\theta)
=
-\langle\theta,x(\theta)\rangle,
\]
we obtain
\[
-\mathcal H(\mathsf n_\theta^*\,\|\,\mathsf m_0)
=
\psi(\theta)+\langle\theta,x(\theta)\rangle
=
\mathcal S(x(\theta)).
\]
This proves equality.
\end{proof}

\begin{rem}
Equality between $\mathcal S_{\mathrm{var}}$ and $\mathcal S$ for every
admissible extensive vector requires an additional duality or
constraint-qualification hypothesis. Without such a hypothesis, the
equality is guaranteed on the equilibrium image
\[
-\nabla\psi(\mathcal D)
\]
but need not hold on the entire boundary of the constraint domain.
\end{rem}
\section{Existence and Uniqueness of Equilibrium Means on Fréchet Manifolds}

%--------------------------------------------------------------
\subsection{Analytic hypotheses and functional framework}
%--------------------------------------------------------------

Let $(\mathcal M,\omega)$ be a weak symplectic Fréchet manifold, let
$\mathcal L$ be the unital vector lattice algebra of observables
introduced in the preceding section, and let $\mathsf m_0$ be a
normalized mean on $\mathcal L$. We fix a possibly unbounded Hamiltonian
$H\in\mathcal L$ and a finite family of possibly unbounded observables
\[
\Phi_1,\ldots,\Phi_N\in\mathcal L,
\]
for instance components of a moment map.

For
\[
(\beta,\lambda)\in\mathbb R\times\mathbb R^N,
\]
set
\[
F_{\beta,\lambda}
=
-\beta H-\sum_{a=1}^{N}\lambda^a\Phi_a.
\]
The admissible thermodynamic domain is
\[
\mathcal D
=
\left\{
(\beta,\lambda)\in(0,\infty)\times\mathbb R^N
\;\middle|\;
F_{\beta,\lambda}\in\mathcal E(\mathsf m_0)
\right\}.
\]

For $(\beta,\lambda)\in\mathcal D$, the central functional is the
\emph{free energy}
\[
\mathcal G_{\beta,\lambda}(\mathsf n)
=
\mathcal H(\mathsf n\,\|\,\mathsf m_0)
+
\beta\,\mathsf n(H)
+
\sum_{a=1}^{N}\lambda^a\mathsf n(\Phi_a),
\]
defined on $\mathsf S(\mathcal L)$. Equivalently,
\[
\mathcal G_{\beta,\lambda}(\mathsf n)
=
\mathcal H(\mathsf n\,\|\,\mathsf m_0)
-
\mathsf n(F_{\beta,\lambda}).
\]

\begin{Definition}[Coercivity and local exponential regularity]
\label{def:coercive}
We say that $(H,\Phi_a,\mathsf m_0)$ satisfies the analytic hypotheses
on a parameter domain if:
\begin{enumerate}
\item[\textup{(i)}] there exists a continuous proper functional
\[
V\colon\mathcal M\longrightarrow[0,\infty)
\]
and constants $c_1,c_2>0$ such that
\[
H(m)\geq c_1V(m)-c_2,
\qquad
m\in\mathcal M;
\]
\item[\textup{(ii)}] the partition functional is locally exponentially
regular on $\mathcal D$ in the sense of
Definition~\ref{def:local-exponential-regularity};
\item[\textup{(iii)}] whenever a compactness argument is used, the
sublevel sets of $\mathcal G_{\beta,\lambda}$ are compact in
$\mathsf S(\mathcal L)$ for the topology of pointwise convergence.
\end{enumerate}
\end{Definition}

\begin{rem}
The lower bound in Definition~\ref{def:coercive}\textup{(i)} is a
model-space coercivity condition. For general finitely additive means,
it does not by itself imply compactness of free-energy sublevel sets.
Such compactness must be verified separately, as required in
condition~\textup{(iii)}.
\end{rem}

%--------------------------------------------------------------
\subsection{Admissible parameters and non-vanishing partition functions}
\label{subsec:admissible-partition}
%--------------------------------------------------------------

For a thermodynamic potential
\[
F_{\beta,\lambda}
=
-\beta H-\sum_{a=1}^{N}\lambda^a\Phi_a,
\]
the positivity of the function $e^{F_{\beta,\lambda}}$ does not, for an
arbitrary normalized mean, automatically imply
\[
\mathsf m_0(e^{F_{\beta,\lambda}})>0.
\]
A positive normalized linear functional need not be faithful and may
vanish on a nonzero positive function.

This issue is particularly important for normalized means defined by
finite-dimensional cut-offs. Let
\[
Z_K(\beta,\lambda)
=
\frac{
\displaystyle
\int_{B_K}
e^{F_{\beta,\lambda}}\,\mathrm d\mu_K
}{
\displaystyle
\int_{B_K}\mathrm d\mu_K
}.
\]
Although
\[
Z_K(\beta,\lambda)>0
\qquad
\text{for every }K,
\]
the limit of the sequence may vanish. Positivity of all the
finite-dimensional partition functionals therefore does not by itself
imply strict positivity of the limiting partition functional.

\begin{Definition}[Admissible thermodynamic domain]
\label{def:admissible-domain}
The admissible thermodynamic domain is
\[
\mathcal D
=
\left\{
(\beta,\lambda)\in(0,\infty)\times\mathbb R^N
\;\middle|\;
F_{\beta,\lambda}\in\mathcal E(\mathsf m_0)
\right\}.
\]
Equivalently, $(\beta,\lambda)\in\mathcal D$ if
\[
e^{F_{\beta,\lambda}}\in\mathcal L,
\qquad
0<
\mathsf m_0(e^{F_{\beta,\lambda}})
<\infty,
\]
and
\[
fe^{F_{\beta,\lambda}}\in\mathcal L
\qquad
\text{for every }f\in\mathcal L.
\]
\end{Definition}

Thus strict positivity of the limiting partition functional is part of
the admissibility condition unless it follows from an independent
structural assumption, such as faithfulness of the reference mean.

\begin{Theorem}[Criteria for finiteness and non-vanishing]
\label{thm:partition-positive}
\label{prop:partition-nonvanishing}
Let $(\beta,\lambda)$ be such that
\[
e^{F_{\beta,\lambda}}\in\mathcal L
\]
and
\[
\mathsf m_0(e^{F_{\beta,\lambda}})<\infty.
\]
Then
\[
0<
\mathsf m_0(e^{F_{\beta,\lambda}})
<\infty
\]
under either of the following assumptions:
\begin{enumerate}
\item[\textup{(i)}] the reference mean $\mathsf m_0$ is faithful on
$\mathcal L$;
\item[\textup{(ii)}] the reference mean is defined by cut-offs, one has
\[
\mathsf m_0(e^{F_{\beta,\lambda}})
=
\lim_{K\to\infty}Z_K(\beta,\lambda),
\]
and
\[
\liminf_{K\to\infty}Z_K(\beta,\lambda)>0.
\]
\end{enumerate}
\end{Theorem}

\begin{proof}
Finiteness is assumed. Since
\[
e^{F_{\beta,\lambda}}>0
\]
pointwise, it is a nonzero positive element of $\mathcal L$. Under
assumption~\textup{(i)}, faithfulness gives
\[
\mathsf m_0(e^{F_{\beta,\lambda}})>0.
\]

Under assumption~\textup{(ii)},
\[
\mathsf m_0(e^{F_{\beta,\lambda}})
=
\lim_{K\to\infty}Z_K(\beta,\lambda)
\geq
\liminf_{K\to\infty}Z_K(\beta,\lambda)>0.
\]
\end{proof}

\begin{rem}
The lower-bound condition in
Theorem~\ref{thm:partition-positive} is sufficient but not necessary. In
applications, it may be replaced by any estimate that directly yields
\[
\mathsf m_0(e^{F_{\beta,\lambda}})>0.
\]
\end{rem}

%------------------------------------------------------------
\subsection{Change of reference under exponential tilting}
\label{subsec:change-reference}
%------------------------------------------------------------


\begin{thebibliography}{99}
\bibitem[AbrahamMarsden78]{AbrahamMarsden78} R.~Abraham and J.~E.~Marsden, \emph{Foundations of Mechanics}, 2nd ed., Benjamin/Cummings Publishing Co., Reading, MA, 1978.
\bibitem[Day57]{Day57} M.~M.~Day, \emph{Amenable semigroups}, Illinois J. Math. \textbf{1} (1957), 509--544.
\bibitem[Hamilton82]{Hamilton82} R.~S.~Hamilton, \emph{The inverse function theorem of Nash and Moser}, Bull. Amer. Math. Soc. (N.S.) \textbf{7} (1982), no.~1, 65--222.
\bibitem[KM97]{KM97} A.~Kriegl and P.~W.~Michor, \emph{The Convenient Setting of Global Analysis}, Mathematical Surveys and Monographs, vol.~53, American Mathematical Society, Providence, RI, 1997.
\bibitem[Magnot17]{Magnot17} J.-P.~Magnot, \emph{The mean value for infinite volume measures, infinite products, and heuristic infinite dimensional Lebesgue measures}, J. Math. \textbf{2017} (2017), Article ID 9853672, 14 p.
\bibitem[MarsdenRatiu99]{MarsdenRatiu99} J.~E.~Marsden and T.~S.~Ratiu, \emph{Introduction to Mechanics and Symmetry}, 2nd ed., Texts in Applied Mathematics, vol.~17, Springer, New York, 1999.
\bibitem[Neeb06]{Neeb06} K.-H.~Neeb, \emph{Towards a Lie theory of locally convex groups}, Japanese J. Math. \textbf{1} (2006), no.~2, 291--468.
\bibitem[Omori97]{Omori97} H.~Omori, \emph{Infinite-Dimensional Lie Groups}, Translations of Mathematical Monographs, vol.~158, American Mathematical Society, Providence, RI, 1997.
\bibitem[Paterson88]{Paterson88} A.~L.~T.~Paterson, \emph{Amenability}, Mathematical Surveys and Monographs, vol.~29, American Mathematical Society, Providence, RI, 1988.
\bibitem[vonNeumann29]{vonNeumann29} J.~von Neumann, \emph{Zur allgemeinen Theorie des Masses}, Fund. Math. \textbf{13} (1929), 73--116.
\end{thebibliography}
\end{document}
